Part~\ref{part:cp2} reproduces the Checkpoint~2 report unchanged, with dated notes where the repository now
shows otherwise. Part~\ref{part:new} reports the work that report does not contain: explicit spiral geometry
and the wind and anchor-removal sweeps of the configurable model (Section~\ref{sec:core}); an A* and a
fixed-action-pattern formulation (Section~\ref{sec:astar_fap}); a benchmark for construction on moving
supports with six execution policies (Section~\ref{sec:mulder}); recovery from execution failures with the
re-entrant planner IPyHOP (Section~\ref{sec:recovery}); and a physical simulation of the planned web under
gravity, wind and support motion (Section~\ref{sec:sim}). The energy advantage of the HTN agent over
undirected search is confirmed at roughly an order of magnitude. Its claimed reliability advantage is
withdrawn, and in the new benchmarks it is closed-loop execution rather than hierarchical structure that
carries construction through disturbance. Negative results are reported as such
(Sections~\ref{sec:record-b} and~\ref{sec:synthesis}).
